Traffic asset inventory is a fundamental task for digital management, risk assessment, and refined maintenance of road infrastructure. For vehicle-mounted panoramic imagery, this task requires both reliable asset discovery and accurate state understanding. However, existing methods still suffer from two major limitations: small, boundary, and densely distributed targets in complex road scenes are prone to missed detection, which weakens inventory completeness; meanwhile, abnormal-state samples are scarce and highly long-tailed, which restricts model performance on low-frequency yet high-risk attributes. To address these issues, this thesis investigates active-perception detection and knowledge-guided attribute recognition enhancement on the WHU-Infra3D dataset, with the aim of improving the completeness, accuracy, and engineering usability of traffic asset inventory.

For the detection task, an active-perception three-stage multi-agent framework is proposed for traffic asset detection. First, equirectangular panoramic images from the Wuhan and Shanghai subsets are converted into four perspective views to alleviate the influence of local geometric distortion in panoramic imagery. Then, through the collaboration between a Detection Agent and a Crop Agent, conventional one-pass full-image detection is reformulated into a pipeline of global initial detection, candidate region recommendation, local refinement, and result fusion, thereby enabling targeted re-inspection of potentially missed regions. In cross-city experiments with 6,712 perspective images from Wuhan for training and 4,700 perspective images from Shanghai for testing, the proposed framework achieves consistent gains on six heterogeneous detectors, namely RexOmni, DINO, Faster R-CNN, DEIM, Grounding DINO, and YOLO26. The average improvements reach 8.18\% in Macro-F1 and 6.41\% in Micro-F1, with more evident advantages in small-object recall, dense-region recovery, and complex-scene detection.

For the attribute recognition task, a domain-knowledge-guided method for long-tail data construction and augmentation is proposed. The method transforms traffic-facility standards, maintenance experience, and visual priors into constrained prompt rules, and performs targeted augmentation for key abnormal states such as corrosion, weathering and aging, structural damage, manhole displacement, and garbage overflow, while preserving the object structure and scene context as much as possible. Based on 7,311 real local samples constructed from WHU-Infra3D, this thesis further generates 4,197 abnormal attribute samples and adopts a unified multimodal supervised fine-tuning paradigm to train the attribute recognition model. Experimental results show that the long-tail enhanced model achieves 88.78\%, 96.12\%, and 93.72\% in instance-level exact-match, attribute-level, and abnormal-attribute accuracy, respectively, improving over the model trained only on real samples by 30.82, 11.60, and 92.71 percentage points. Meanwhile, the performance gap between normal and abnormal states is reduced significantly from 95.63\% to 6.08\%.

In summary, this thesis establishes a composite perception framework for traffic asset inventory from two aspects, namely active-perception detection and knowledge-guided long-tail enhancement. The results demonstrate that the proposed methods can effectively improve asset discovery completeness, abnormal attribute recognition, and structured record quality while maintaining engineering feasibility, thereby providing methodological support and data foundations for digital traffic asset inventory, intelligent inspection, and maintenance decision-making.
